% Source: https://artofproblemsolving.com/community/c3h2389939_alternate_way_to_prove_the_quadratic_formula
% Stable identity: aops-topic-2389939-q01
A student proves the quadratic formula as follows. Let $ax^2+bx+c=0$ with $a\ne0$ have roots $x_1,x_2$. By Vieta's formulas, $x_1+x_2=-\dfrac ba$ and $x_1x_2=\dfrac ca$, so
\[(x_1-x_2)^2=(x_1+x_2)^2-4x_1x_2=\dfrac{b^2-4ac}{a^2}.\]
Taking square roots, $x_1-x_2=\dfrac{\sqrt{b^2-4ac}}{a}$ (the sign may be chosen by naming the roots suitably). Adding this to $x_1+x_2=-\dfrac ba$ and subtracting gives $x_{1,2}=\dfrac{-b\pm\sqrt{b^2-4ac}}{2a}$.

\begin{enumerate}
\item[(a)] What does this argument assume without proof? Explain why the assumption matters, using the equation $x^2+1=0$ over the real numbers.
\item[(b)] Complete the argument into a proof of the following statement: if $a\ne0$, $b$, $c$ are real and $b^2-4ac\ge0$, then the real solutions of $ax^2+bx+c=0$ are exactly the numbers $\dfrac{-b\pm\sqrt{b^2-4ac}}{2a}$.
\end{enumerate}
